\documentclass[aps,prb,onecolumn,nofootinbib,superscriptaddress]{revtex4-2}
\usepackage{amsmath,amssymb,amsthm,mathtools,bm}
\usepackage{booktabs}
\usepackage[colorlinks=true,linkcolor=blue,citecolor=blue,urlcolor=blue]{hyperref}
\usepackage{microtype}

\newcommand{\repo}[1]{\texttt{\detokenize{#1}}}

\begin{document}

\title{Worker 15 Manifest and Data Audit}
\author{Worker 15}
\date{June 11, 2026}

\begin{abstract}
I audited manifest-level and tabular metadata for non-excluded cache, CPU-data, GPU-data,
and campaign outputs, without reading binary arrays or executing notebooks.  The main
conclusion is that the newer campaign products are the most reliable primary evidence:
they usually record the canonical CPU/GPU dynamics entry point, run configuration, sample
counts, saved-observable policy, and file-size accounting.  By contrast, the older
\repo{cache/G_history_samples} layer is extremely large, filename-driven, weakly
manifested, and contains multiple rerun/duplicate/test-style artifacts.  Several
metadata-visible caveats should be carried into any synthesis: postselection runs often
have requested samples greater than actual samples because the canonical driver stores a
single deterministic sample; at least one analysis manifest records no runs despite
derived comparison outputs; and one CPU convergence manifest was produced from a dirty
working tree.
\end{abstract}

\maketitle
\setcounter{tocdepth}{2}
\tableofcontents

\section{Source Scope}

The audit followed \repo{REPO_POLICY.md} before inspection.  I excluded every path
containing \repo{haining}, \repo{haoyu}, or \repo{erroneous}, case-insensitive, and did
not inspect \repo{repo_synthesis/main_pass/worker_notes},
\repo{repo_synthesis/main_pass/supervisor_notes}, or \repo{repo_synthesis/main_pass/boss}.
No files were edited, no notebooks were executed, and binary arrays were treated only by
filename, byte size, and text manifests.

The inspected evidence consists of JSON manifests, CSV inventories, run indices, run
summaries, and file-size listings under \repo{cache/}, campaign-specific \repo{cpu_data}
and \repo{gpu_data} trees, \repo{dw_convergence}, \repo{colab_charge_fluctuations},
\repo{colab_small_system_testing}, and
\repo{colab_large_entanglement_scaling_N20}.  A manifest inventory found 30
\repo{campaign_manifest.json} files and 333 \repo{run_summary.json} files in the principal
campaign trees.  The legacy \repo{cache/G_history_samples} binary layer contains 38
non-excluded \repo{.npz}/\repo{.npy} files larger than 1 GiB, totaling approximately
662.3 GiB by filename/size metadata alone.

\section{Findings}

\subsection{Primary Data Layer}

The strongest primary layer is the newer campaign/run-summary structure.  Representative
GPU manifests explicitly identify
\repo{classA_U1FGTN_gpu.run_markov_circuit} as the canonical dynamics entry point, record
A100/complex128 batching policy, and set \repo{covariance_history_saved=false} for
streaming observable campaigns.  The large-entanglement run summary records the canonical
GPU entry point in its configuration, the A100 40 GB strategy, the final-only Markov
result, and complete status.  The charge-fluctuation GPU campaign records six
multi-geometry configurations, saved observables, per-run byte counts, and a total saved
observable footprint of 163,681,202 bytes.

The strongest CPU primary layer is the structured
\repo{purification_charge_sharpening_alpha_sweep} campaign.  Its campaign manifest records
the canonical CPU entry point \repo{classA_U1FGTN.run_markov_circuit}, declares the run
complete, suppresses covariance-history saving, and documents the active domain-wall slab
observer.  The companion CSV inventory gives a machine-checkable row-level run index with
columns for protocol, geometry, alpha, cycle count, expected/completed samples, and
completion.

\subsection{Corroborating and Derived Data}

Analysis-output manifests should be treated as corroborating rather than primary unless
they point cleanly back to a primary campaign.  The streaming covariance characterization
manifest is useful because it names the source dataset, lists all derived figures/tables,
records analysis targets, and gives a six-record map from derived scaling fits back to
GPU \repo{local_charge_cell.npz} inputs.  It also records a skipped \repo{N20x60} run
because both \repo{local_charge_cell.npz} and
\repo{streaming_covariance_scaling_fits.npz} were missing; this is valuable negative
evidence, not a completed datum.

The \repo{dw_convergence} analysis manifest is internally informative but lower priority
as primary evidence.  It records the canonical CPU engine and a clear N20x20 convergence
configuration, but its \repo{runs} array is empty despite nearby derived CSV outputs, and
its git status contains many modified, deleted, and untracked paths.  That combination
makes it useful for method provenance and for locating derived products, but not as a
self-contained completed campaign manifest.

\subsection{Redundancy and Stale-Legacy Risk}

The legacy \repo{cache/G_history_samples} layer is the highest-risk stratum for stale or
redundant evidence.  The file-size listing alone shows many tens-of-GiB covariance-history
artifacts, including names containing \repo{slope_drift_testing}, \repo{fromcache-rerun},
\repo{legacydup1}, and paired variants such as \repo{init-default} versus
\repo{init-rerun}.  These names encode useful hints, but they are not a substitute for
structured provenance.  Such files should not be used as primary evidence unless a
current manifest explicitly points to them and resolves the intended run identity.

The newer campaign layer also contains deliberate redundancy: run summaries, campaign
manifests, run indices, and analysis inventories repeat configuration and sample-count
fields.  This redundancy is beneficial when consistent, because it allows cross-checks
without opening arrays.  It becomes a warning signal when sample counts differ or when
analysis records list skipped inputs.

\subsection{Metadata-Visible Contradictions}

The clearest recurring contradiction is requested versus actual sample count for
postselection.  The purification-dynamics GPU campaign requests 100 samples for
postselection but states that the canonical GPU driver forces \repo{samples=1}; the result
rows then show \repo{samples_requested=100} and \repo{samples_actual=1}.  The covariance
history campaign has the analogous warning for 10 requested postselection samples saved as
one deterministic sample.  These are not necessarily physical contradictions, but they are
semantic traps: postselection rows should not be ensemble-weighted as if they had the same
sample count as perfect-correction rows.

\section{Evidence Table}

\begin{table}[h]
\caption{Representative manifest and metadata evidence.}
\begin{tabular}{p{0.36\textwidth}p{0.17\textwidth}p{0.39\textwidth}}
\toprule
Source & Lines & Audit use \\
\midrule
\repo{REPO_POLICY.md} & 1--19 & Establishes canonical CPU/GPU entry points and requires production metadata to record the entry point. \\
\repo{colab_charge_fluctuations/gpu_data/index.json} & 1--10 & Top-level GPU index points only to \repo{purification_dynamics_maxmix/latest_campaign.json}; useful as a bundle pointer but incomplete as a full inventory. \\
\repo{colab_small_system_testing/gpu_data/index.json} & 1--10 & Similar top-level pointer for \repo{streaming_covariance_protocol_characterization}; not a comprehensive data catalogue. \\
\repo{colab_charge_fluctuations/gpu_data/purification_dynamics_maxmix/.../campaign_manifest.json} & 11, 92--100, 104--105, 270--364 & Records canonical GPU engine, no covariance-history saving, postselection sample downgrade note, six results, per-run byte counts, saved observables, and 163,681,202 total saved bytes. \\
\repo{colab_charge_fluctuations/gpu_data/streaming_covariance_observables/.../N20x50.../run_summary.json} & 13--33 & Records canonical GPU engine, no covariance-history saving, 10,000 scalar metric rows, and streaming observer mode. \\
\repo{colab_charge_fluctuations/analysis_outputs/streaming_covariance_characterization/.../analysis_manifest.json} & 36--106, 107--133 & Maps six derived analysis records back to GPU inputs and records a skipped N20x60 run with missing files. \\
\repo{colab_charge_fluctuations/cpu_data/purification_charge_sharpening_alpha_sweep/.../campaign_manifest.json} & 25--35, 36--49 & Records complete CPU campaign, canonical CPU engine, no covariance-history saving, active-slab observer, and completed run rows. \\
\repo{colab_charge_fluctuations/.../tables/inventory.csv} & 1--35 & CSV inventory confirms row-level alpha/geometry/cycle/sample completion for the fine CPU sweep. \\
\repo{colab_small_system_testing/gpu_data/covariance_protocol_histories/.../campaign_manifest.json} & 8, 217--219, 266--299, 343--377 & Records canonical GPU engine, full history storage for selected runs, explicit shard sizes/shapes, and postselection reduction from requested 10 samples to 1 saved sample. \\
\repo{colab_large_entanglement_scaling_N20/.../run_summary.json} & 49--67, 196--233 & Records checkpoint/output paths, canonical GPU engine, A100 40 GB batching, final-only Markov result metadata, and complete status. \\
\repo{dw_convergence/n20x20_maxmix_dwtrunc_convergence_cpu/analysis_manifest.json} & 2--24, 76--90 & Records canonical CPU engine and configuration, but also dirty git status and an empty \repo{runs} array. \\
\repo{cache/G_history_samples} size listing & manifest-only file sizes & Shows 38 non-excluded binary files larger than 1 GiB, totaling about 662.3 GiB; filenames include rerun, testing, and duplicate-style labels. \\
\bottomrule
\end{tabular}
\end{table}

\section{Interpretation}

For synthesis, use the structured campaign manifests and run summaries as the primary
data spine.  Prefer campaigns that satisfy all of the following: canonical dynamics entry
point present; sample counts explicit; covariance-history policy explicit; run status or
completion explicit; and derived analysis records linked by relative paths to primary
run directories.  The charge-fluctuation GPU streaming/purification campaigns, the CPU
alpha sweep, and the large-entanglement N20 run meet this standard most clearly.

Use analysis manifests and tables as corroborating summaries.  They are especially useful
when they name source datasets and missing inputs, but they should not override primary
run manifests on sample count or protocol identity.  Treat legacy \repo{cache} binaries as
archival or diagnostic unless referenced by current manifests.  Their scale and filename
semantics imply substantial duplication and historical experimentation, and their lack of
structured provenance makes them poor foundations for a final narrative.

\section{Caveats}

This audit did not open \repo{.npz}, \repo{.npy}, \repo{.parquet}, images, videos, PDFs,
or notebooks.  Therefore it verifies provenance, shape/size claims where written in
manifests, sample-count semantics, and manifest consistency, but not numerical values
inside binary arrays.  Line references are to the inspected text files as present on
June 11, 2026.  Some top-level indices are intentionally sparse and should not be
mistaken for exhaustive inventories.

\section{Confidence}

Confidence is high for metadata-level conclusions: canonical engine recording, sample
count mismatches, missing-input records, dirty git provenance, and binary-size risk are
directly visible from text manifests and file-size listings.  Confidence is moderate for
ranking data products as primary versus corroborating, because that judgment depends on
provenance completeness rather than numerical validation.  Confidence is low for any
claim about the numerical correctness of observables, since binary arrays and notebooks
were intentionally not inspected.

\end{document}
